\scanpage{216}
\mspage{\deleted{\unsure{7|8}}71\quad A}

$B$ est \uline{$R$-clos si} $B=B^\flat$. On aura pour tout \deleted{$A$} $A\subset P$, $B\subset Q$

\[
\left\{\begin{aligned}
 (A^{\unsure{\cup|\smile}})^\flat&=A^\cup\quad\text{i.e. }A^\cup\text{ est }R\text{-clos}\\
 (B^{\unsure{\cap|\frown}})^{\unsure{\sharp}}&=B^\cap\quad\text{i.e. }B^\cap\text{ est }R\text{-clos.}
\end{aligned}\right.
\tag*{(5.6)}
\]

En particulier

\[
 (A^\sharp)^\sharp=A^\sharp,\qquad (B^\flat)^\flat=B^\flat
\tag*{(5.7.)}
\]

i.e. la $R$-\uline{clôture} $A^\sharp$ de $A$ \inserted{$\subset P$,} est $R$-close, et de $=$ la $R$-clôture $B^\flat$ de $B\subset Q$.

En fait, je m'aperçois que je suis en train de mélanger deux espèces de faits différents. \deleted{\unsure{D'}} Les relations (5.2.), et (\unsure{5.4.|5.3.}) à (5.7), ne dépendent pas de la donnée des relations d'ordre sur $P,Q$, ni d'axiomes reliant celles-ci à la relation donnée $R$, mais seulement de cette relation, laquelle n'a pas à être soumise à aucune condition. Notons que

\[
 \varnothing^\cup=\deleted{\unsure{B}}Q,\qquad \varnothing^\cap=\deleted{\unsure{A}}P
 \tag*{(5.8.)}
\]

tandis que $\deleted{\unsure{\varnothing^\sharp=}}$

\[
 \varnothing^\sharp=\unsure{Q}^\cap,\qquad \varnothing^\flat=P^\cup
 \tag*{(5.9)}
\]

sont respectivement la plus petite partie $\deleted{\unsure{cl}}R$-close de $P$ et de $Q$.

L'application

\[
 A\unsure{\vdash|\mapsto} A^\sharp
 \tag*{(5.10\unsure{.|:})}
\]

est une application \uline{croissante} \inserted{(et \uline{idempotente})} de $\deleted{\unsure{\mathfrak P A}}$ \deleted{dans} $\unsure{\mathfrak P|\mathcal P}(P)$ dans lui-même, ayant \inserted{\unsure{donc|aussi}} la propriété

\deleted{\[
 \unsure{(A\subset B)^\sharp=A}
\]}
% archive sheet number: 216
